% Transcription — fonds Alexandre Grothendieck, shelfmark 84, batch 3 (pages 41-60).
% Established from the facsimile archives/batches/84/batch-03.pdf, per the skill
% transcribe-grothendieck. Page 1 of the batch file is the archivists' page 41;
% their pencil numbers bottom-left (« 41 » ... « 60 ») agree throughout.
%
% Pass: Opus 5.5 (claude-opus-5-5), 2026-10-09 — first pass, unchecked against the pages by a human.
%
% Read from the batch PDF rendered at 220 dpi and cropped by half-sheets
% (the 700-dpi tiles were not yet available when the pass was made).
% Absent: page 44 (blank sheet, only show-through). Page 42 is a cover
% sheet in his hand on music paper: it carries no \page{}, and its words
% open the section heading, as hand.md prescribes.
% Several + / - signs in pages 43-53 are blotted over with ink; they are
% given as the computation requires, and a note says so where it first
% happens. The faint rotated writing in the left margin of page 53 is
% show-through of page 54 and is not transcribed there.

\documentclass[11pt,a4paper]{article}
\input{../preamble/grothendieck}

\folder{84}
\batch{3}
\pages{41}{60}
\foldertitle{[Formes quadratiques] : notes manuscrites (s.d.).}
\dating{[à partir de 1982-vers 1986]}
\watermark{Édition de démonstration}

\begin{document}

\page{41}
\note{la page continue un argument commencé avant ce lot : $f$, $P^{*}$, $A^{*}$, $W(A)^{*}$ et $SW(A)$ y sont introduits.}

et on a
\[
  f(u.D) = u\bar{u}^{-1} f(D) \qquad D \in \Gamma P^{*},\ u \in \Gamma A^{*}
\]
(\,$= f(D)$ \ill{}\,), i.e.\ le morphisme de schémas est compatible avec les opérations des groupes d'op.\ $W(A)^{*}$ et $SW(A)$ et l'homomorphisme de groupes
\[
  W(A)^{*} \longrightarrow SW(A), \qquad u \longmapsto u\bar{u}^{-1}
\]
), donc aussi (en restreignant le groupe d'opérateurs sur $P^{*}$ à $SW(A)$) …
\[
  u \longmapsto u^{2}, \qquad SW(A) \longrightarrow SW(A)
\]

\section{$X$-scindages, produits contractés d'extensions $X$-scindées, produits contractés de revêtements quadratiques, invariants d'Hasse-Stiefel-Whitney généralisés, etc.}
\note{ce titre est la seule inscription d'une feuille de papier à musique, de sa main, qui précède la page 43 (page 42 du fonds) ; il y écrit d'abord « \struck{produits} contractés de revêtements quadratiques », avec le premier mot biffé puis récrit. La feuille porte aussi au crayon le numéro « (79) » entre parenthèses.}

\subsection{Revêtements quadratiques}

\page{43}
\[
  0 \to \mathcal{O} \to A \xrightarrow{\ \Psi\ } \check{L} \to 0
  \qquad E = \check{A}
  \qquad 0 \to L \to E \xrightarrow{\ \psi\ } \mathcal{O} \to 0
\]
\note{sous $\check{L}$ : « \uncertain{inv} » ; la lettre au-dessus de la première flèche est \uncertain{$\Psi$}.}
\[
  \begin{array}{ll}
    P = \psi^{-1}(1) \quad \text{tors.\ sous } L \qquad & \psi(x) = \langle 1, x\rangle = x(1)\\[2pt]
    P_{\lambda} = \psi^{-1}(\lambda) \quad \text{id.} &
  \end{array}
\]
\begin{itemize}
  \item $T \in \Gamma P_{2} \subset E$ \quad fonction trace
  \item $N \in \Gamma \mathrm{Sym}^{2}(E) \subset \mathrm{Sym}^{\bullet}(E)$ (fonctions pol.\ sur $A$) \quad forme quadratique \emph{norme}
\end{itemize}

N.B.\ 1) $0 \to E \otimes L \to \mathrm{Sym}^{2}(E) \xrightarrow{\ \chi\ } \mathcal{O} \to 0$ \quad où $\chi(Q) = Q(1)$

on a $N \in \chi^{-1}(1)$ (tors.\ sous $E \otimes L$) i.e.\ $N(1) = 1$.

2) \struck{\ill{}} On a $0 \to L^{\otimes 2} \to E \otimes L \to L \to 0$ et $\mathrm{Sym}^{2}(E)$ a suite de comp.\ à quotients successifs $\mathcal{O}$, $L$, $L^{\otimes 2}$. On identifie $L^{\otimes 2}$ à ss-module de $\mathrm{Sym}^{2}E$.

\marginal{\emph{attention au signe} (encadré) — \uncertain{Sur} la discr.\ \ill{} changé du signe}
\begin{itemize}
  \item $-\delta \in \Gamma(L^{\otimes 2})$ \quad « discriminant » de $N$
\end{itemize}

N.B.\ Forme \uncertain{bilinéaire} sym.\ $\varphi_{N}$ associée à $N$
\[
  \varphi_{N} : A \times A \to \mathcal{O} \quad\text{d'où}\quad A \to \check{A} \overset{\mathrm{df}}{=} E
\]
d'où
\[
  \textstyle\bigwedge^{2} A \to \bigwedge^{2} E
  \qquad\text{d'où}\qquad \delta \in \Gamma L^{\otimes 2}
\]
\[
  \textstyle\bigwedge^{2} A \simeq \check{L}, \qquad \bigwedge^{2} E \simeq L,
  \qquad L^{\otimes 2} \simeq \mathrm{Hom}(\check{L}, L)
\]
\[
  \boxed{T^{2} - 4N = +\delta} \qquad\text{i.e.}\qquad \boxed{4N = T^{2} - \delta}
\]
\marginal{« attention signe ! », avec une flèche vers le $+$ de la première formule.}
\note{ici et dans la suite du lot, plusieurs signes $+$ ou $-$ sont recouverts d'une tache d'encre (surcharge) ; on les donne tels que le calcul les exige, sans les avoir lus. C'est le cas du signe de la seconde formule encadrée.}
\marginal{N.B.\ on a $\varphi_{N}(x,y) =$ \uncertain{$x\bar{y} + y\bar{x}$}, $\mathrm{Tr}(x)\mathrm{Tr}(y) = \varphi_{N}(x,y) + \mathrm{Tr}(xy)$ \ill{}, et on a $\varphi_{N}(x,y) = \mathrm{Tr}(x)\mathrm{Tr}(y) - \mathrm{Tr}(xy)$ \uncertain{et on a}.}
\note{note marginale écrite en diagonale ; lecture d'ensemble douteuse.}

\emph{N.B.} Soient $x', x'' \in \Gamma E$ \struck{\ill{}} deux sections \emph{distinctes} en tt pt, et qui correspondent à des homs d'alg.\ $A \overset{x',x''}{\rightrightarrows} \mathcal{O}_{X}$ (donc $x', x'' \in \Gamma P_{1}$). On a donc $A \simeq \mathcal{O}_{X} \times \mathcal{O}_{X}$
\[
  \text{\struck{$x'(a,b)$}} \qquad x'(a,b) = a, \quad x''(a,b) = b,
  \qquad \mathrm{Tr}(a,b) = a + b, \quad N(a,b) = ab
\]
d'où
\[
  \boxed{\begin{array}{l} T = x' + x'' \\ N = x'x'' \\ \delta = (x' - x'')^{2} \end{array}}
  \qquad \text{\struck{\ill{}}}
\]

\page{45}
Considérons la fonction polynomiale sur $P_{1}$
\[
  x \longmapsto \underbrace{(x - x')}_{\in L}\,\underbrace{(x - x'')}_{\in L} \in L^{\otimes 2}
\]
i.e.
\[
  x \longmapsto x^{2} - Tx + N \qquad \text{\struck{\ill{}}}
\]
restriction de la forme quadratique
\[
  Q : x \longmapsto x^{2} - \psi(x)Tx + \psi(x)^{2}N, \qquad E \to L^{\otimes 2} \subset \mathrm{Sym}^{2}(E)
\]
\note{au-dessus de la formule, deux essais biffés, \struck{\ill{}} et \struck{\ill{}}, avant $\psi(x)$.}
sur $E$. Cette forme est définie \add{via cette formule} (en tous cas, \uncertain{sans} supposer existence de $x'$, $x''$), elle prend ses valeurs dans $\mathrm{Sym}^{2}(E)$, mais \struck{on} on vérifie qu'elle est à valeurs dans $L^{\otimes 2} \subset \mathrm{Sym}^{2}$ (argument « universel » …). Sa restriction à $L$ est
\[
  x \longmapsto Q(x) = x^{2} : L \to L^{\otimes 2}
\]
(partie « homog.\ » de la fonction polynomiale \add{de degré 2} \uncertain{canonique} sur $P_{1}$).

\emph{N.B.} Équiv.\ de cat.\ fibrées
\[
  \left\{\begin{array}{l}\text{rev.\ quadratiques}\\ \text{de } S\end{array}\right\}
  \;\simeq\;
  \left\{\begin{array}{l}
    \text{\struck{Triples} $(\text{\struck{$P_{1}$}}\ L, P_{1}, Q)$ sur $S$}\\
    \text{a) $L$ Module inv.}\\
    \text{b) $P_{1}$ $L$-torseur}\\
    \text{c) $Q : P_{1} \to L^{\otimes 2}$ fonct.\ pol.\ de degré 2}\\
    \qquad\text{à partie homogène $x \mapsto x^{2} : L \to L^{\otimes 2}$}
  \end{array}\right.
\]
(\,$Q$ est section d'un torseur sous \struck{$A$} $\mathrm{Hom}(E, L^{\otimes 2})$ [$E$ : module enveloppant de $P_{1}$] $\simeq \check{E} \otimes L^{\otimes 2}$ $\simeq$ \ill{} $A \otimes L^{\otimes 2}$, ext.\ de $L$ par $L^{\otimes 2}$ $\ast$ … ). Le rev.\ quadratique associé à une donnée peut être interprété comme le « schéma des zéros de $Q$ », \uncertain{plongé lui-même comme sous-schéma} du fibré affine $\mathbb{P}_{1}$ correspondant à $P_{1}$.

\marginal{$\ast$ Si $T$ ($= x' + x''$) est \uncertain{donné} (dans $\Gamma P_{2}$), \ill{} $N$ \ill{} $\delta \in \Gamma(L^{\otimes 2})$ \ill{} $T^{2} - 4N = \delta$ \ill{} $x^{2} - xT + N$ \ill{} Mais \ill{}}
\note{note écrite en diagonale dans le coin inférieur gauche, en grande partie illisible.}

\page{46}
Re N.B. Considérons
\[
  \varphi_{N} : A \longrightarrow \check{A} = E
  \qquad\text{via forme bil.\ ass.\ à forme qu.\ $N$}
\]
d'où
\[
  A \xrightarrow{\ \varphi_{N}\ } E \xrightarrow{\ Q\ } L^{\otimes 2}
\]
On trouve
\[
  \boxed{Q(\varphi_{N}(x)) = - N(x)\,\delta}
\]
\marginal{vérification à partir du cas décomposé,}
\note{le signe devant $N(x)\delta$ est recouvert ; p.~52 il écrit $+N(x)\delta$, alors que p.~46 plus bas « $Q$ correspond à $-N$ ».}

\struck{Si $\delta$ inv.\ (cas étale)}

Conditions équivalentes
\begin{enumerate}
  \item[a)] $X'$ rev.\ \emph{étale} de $X$
  \item[b)] $\varphi_{N} : A \to \check{A} = E$ \emph{iso}
  \item[c)] $\delta \in \Gamma L^{\otimes 2}$ section \emph{inv.}, i.e.\ ne s'annulant nulle part.
\end{enumerate}
Alors $\delta$ base de $L^{\otimes 2}$, $\varphi_{N} : A \simeq E$, et \add{\uncertain{vu}} si on fait les identifications, $Q$ correspond à $-N$.

Si par contre $\delta = 0$ (cas tot\supplied{alemen}t ramifié) on trouve
\[
  Q(\varphi_{N}(x)) = 0 \,..
\]

Considérons
\[
  x \longmapsto \boxed{2x - T = z} \quad (z \in \Gamma L), \qquad P_{1} \longrightarrow L
\]
c'est un iso si 2 inv.\ sur $X$ donc iso sur $X_{2} = X - V(2)$. Si $x$ et $z$ se correspondent, on a
\[
  2x = z + T \;\Longrightarrow\; Q(z + T) = 4Q(x)
\]
et \add{si $x$ solution de} l'éq.\ quadratique $X'$ de $X$ i.e.\ $Q(x) = 0$, on a $Q(T + z) = 0$. On en trouve, \struck{puisque} (comme $\psi(z + T) = 2$)
\[
  \text{\struck{$Q(T+z)$}} \qquad Q(z + T) = z^{2} - \delta
\]
\[
  z^{2} + 2Tz + T^{2} - 2(Tz + T^{2}) + 4N = z^{2} - (T^{2} - 4N) = z^{2} - \delta
\]
donc
\[
  \left\{\begin{array}{l}
    Q(x) = 0 \ \text{ i.e.\ } x \in \Gamma(X'/X) \;\Longrightarrow\; \boxed{z^{2} = \delta}\\
    \text{(et iso si 2 inv.\ sur $X$)}
  \end{array}\right.
\]

\page{47}
Le groupe $\mathbb{Z}/2\mathbb{Z}$ opère sur $X'/X$, via
\[
  \left\{\begin{array}{ll}
    x \longmapsto \text{\struck{\ill{}}}\ \bar{x} = T - x & \text{sur } P_{1} \quad (\text{auto})\\
    x + \bar{x} = T & \text{si } x \in \Gamma P_{1}
  \end{array}\right.
\]
on a encore
\[
  \left\{\begin{array}{ll}
    x \longmapsto \bar{x} = \psi(x)T - x & \text{sur } E\\
    \text{i.e.\ } x + \bar{x} = \psi(x)T &
  \end{array}\right.
\]
On a
\[
  \left\{\begin{array}{l}
    \psi(T) = T\\
    \psi(N) = N\\
    Q(\text{\struck{$\psi(x)$}}\ \bar{x}) = Q(x)
  \end{array}\right.
\]
\note{la lettre lue $\psi$ dans les deux premières lignes est douteuse ; une flèche revient sur l'accolade.}
Sur $A$, l'involution donnée par
\[
  x \longmapsto \bar{x} = \mathrm{Tr}(x) - x \qquad\text{i.e.}\qquad x + \bar{x} = \mathrm{Tr}(x).1
\]
\note{sous le « $1$ » est écrit un signe qui ressemble à $\delta_{0}$.}
Considérons $X_{0} = V(2)$, \struck{\ill{} $N(x) \in \Gamma L$ et $T_{0}, N_{0}$,} et $E_{0}, L_{0}, T_{0}, N_{0}, Q_{0}$ les restrictions à $X_{0}$. On a \struck{\ill{} $T_{0} \in E_{0}$} dans
\[
  T_{0} \in \Gamma L_{0} \quad (\Gamma E_{0})
\]
et
\[
  T_{0}^{2} = \delta_{0} \ (= -\delta_{0}) \quad\text{sur } X_{0}
\]
et de même : si $T_{0}$ est relevé loc\supplied{alemen}t \add{en relèvement $\tau_{1}$ de $E$} (sur $X_{1} = V(4)$), \struck{\ill{}}
\[
  (\text{p.\ ex.\ } \tau_{1} = T|X_{1} \overset{\mathrm{df}}{=} T_{1})
  \qquad\text{on a}\qquad
  \tau_{1}^{2} = \delta_{1} \quad (\overset{\mathrm{df}}{=} \delta|X_{1})
\]
(N.B.\ cela ne dépend pas du choix de ce relèvement, car $(\tau_{1} + 2\alpha)^{2} \equiv \tau_{1}^{2}$ (4)\,) On l'exprimera par l'écriture
\[
  \boxed{T_{0}^{2} \equiv \delta \quad (4)}
\]

Ici posé, si $x$ section de $X'/X$, $z$ la section corr.\ de $L$, tq $z = 2x - T$, \struck{$z_{0}$} $z_{0}$ la restr.\ à $X_{0}$, on aura
\[
  \left\{\begin{array}{l} z_{0} = T_{0} \\ z^{2} = \delta \end{array}\right.
\]
i.e.\ $z$ est une solution de $z^{2} = \delta$ \add{sur $X$} qui prolonge la solution déjà donnée $T_{0}$ sur $X_{0}$.

\page{48}
Soit
\[
  \mathfrak{X}' : \mathrm{Sch}_{/X} \to \mathrm{Ens}
\]
dont la val.\ sur $X$ soit
\[
  \mathfrak{X}'(X) = \Bigl\{ x \in \Gamma L \Bigm| \begin{array}{l} x^{2} = \delta \\ x|X_{0} = T_{0} \end{array} \Bigr\}
\]
et idem pour $Y/X$.
\[
  \mu'_{2} \subset \mu_{2,X} : \mathrm{Sch}_{/X} \to \mathrm{Groupes}, \quad\text{comme}
\]
\[
  \mu'_{2}(X) = \bigl\{ \lambda \in \Gamma \mathcal{O}_{X}^{*} \bigm| \lambda^{2} = 1,\ \lambda_{0} = 1 \bigr\}
\]
Alors $\mu'_{2}$ opère sur $\mathfrak{X}'$, et si $\delta$ est inv.\ i.e.\ si $X'/X$ étale, donc $\mathfrak{X}'$ est un torseur sous $\mu'_{2}$ (alors que $X'/X$ est un torseur sous $(\mathbb{Z}/2\mathbb{Z})_{X}$). On a compatibilité
\[
  X' \xrightarrow{\ f\ } \mathfrak{X}', \qquad (\mathbb{Z}/2\mathbb{Z})_{X} \longrightarrow \mu'_{2X}
\]
i.e.
\[
  f(\bar{x}) = -f(x) \qquad \text{\struck{\ill{}}}
\]
Ici
\[
  z \longmapsto \bar{z} \ : \ \mathfrak{X}' \to \mathfrak{X}'
\]
est l'automorphisme induit par la section $-1$ de $\mu'_{2X}$ \struck{\ill{}}.
\marginal{N.B.\ $\mathfrak{X}'$ et $\mu'_{2}$ \ill{} \uncertain{fpqc} \ill{} $2 \in \mathfrak{m}_{X,x}$, $2 \neq 0$}
\note{note écrite verticalement dans la marge gauche, en grande partie illisible.}

Supposons 2 $\mathcal{O}_{X}$-régulier. Alors la donnée d'un rev.\ quadratique $X'/X$ \emph{équivaut} à celle de
\[
  \left\{\begin{array}{l}
    L,\ \delta \in \Gamma L^{\otimes 2},\ T_{0} \in \Gamma(L_{0}) \quad (\text{où } L_{0} = L|X_{0})\\
    \text{tels que, satisfaisant } T_{0}^{2} \equiv \delta \ (4)
  \end{array}\right.
\]
On va décrire les sections $x$ de $X'/X$ en termes de la section $z$ de $L$, dans le cas où $X' \to \mathfrak{X}'$ \uncertain{pas iso} i.e.\ si $X_{0} \neq \emptyset$. \struck{Bien sûr} Bien sûr \ill{}

\page{49}
\[
  \Gamma(X'/X) \hookrightarrow \Gamma(\mathfrak{X}'/X) \qquad (\text{injectif})
\]
car $X - X_{0}$ schém.\ dense dans $X$. Donc $x$ connue quand on connaît $z$. Quelles conditions imposer sur $z$ ? \struck{Il suffit que} Notons que
\[
  \text{\struck{\ill{}}} \qquad T + z \equiv 0 \quad (2)
\]
car cela signifie $T_{0} \pm z_{0} = 0$ i.e.\ $z_{0} = T_{0}$ (OK), donc on pourra prendre
\[
  x = \frac{T + z}{2},
\]
je dis que a) c'est une section de $P_{1}$, i.e.\ $\psi(x) = 1$ (or cela équivaut à : $2\psi(x) = 2$ i.e.\ $\psi(2x) = 2$, OK, $\psi(2x) = \psi(T + z)$), et b) on a bien $Q(x) = 0$ (or cela équivaut à : $4Q(x) = 0$ i.e.\ $Q(2x) = 0$ i.e.\ $Q(T + z) = 0$, donc à $z^{2} = \delta$ … ).

Description de $E$ en termes de $L$, $\delta$, $T_{0}$ : soit $\tau$ une section (loc\supplied{alemen}t) de $L$ qui remonte $T_{0}$
\[
  \Gamma E = \Bigl\{ (\lambda, \xi) \Bigm| \lambda \in \Gamma \mathcal{O},\ \xi \in \Gamma\Bigl(\frac{\lambda\tau}{2} + L\Bigr) \Bigr\}
  \qquad
  \left\{\begin{array}{l} \lambda = \psi(x) \\ \xi = x - \lambda\frac{T}{2} \end{array}\right.
\]
\marginal{sous $\frac{\lambda\tau}{2} + L$ : « ne dépend pas du choix de $L$ » \note{sic ; on attend « de $\tau$ ».} — N.B.\ $\frac{\lambda\tau}{2} + L \subset \frac{1}{2}L$}

on a aussi
\[
  \boxed{\Gamma E \simeq \bigl\{ (\lambda, \zeta) \bigm| \lambda \in \Gamma \mathcal{O},\ \zeta \in \Gamma(\lambda\tau + 2L) \subset \Gamma(L) \bigr\}}\ {}^{\ast}
\]
\[
  \boxed{\left\{\begin{array}{l} \lambda = \psi(x) \\ \zeta = 2x - \psi(x)T \end{array}\right.}
\]
$(\mathrm{d\acute{e}f})$ $L_{\lambda}$ sous-\uncertain{faisceau} de $L$, i.e.\ sections de $L$ qui remontent la section $\lambda T_{0}$ de $L_{0}$, i.e.\ $\boxed{\zeta_{0} = \lambda_{0}T_{0}}$.
\marginal{$\ast$ via l'iso.\ de $L$ \uncertain{dessus} \ill{} $x = x + 2\xi$ \ill{}}

Relation entre $(\lambda, \xi)$ et $(\lambda, \zeta)$ : $\zeta = 2\xi$.
\[
  0 \to L \longrightarrow E \xrightarrow{\ \psi\ } \mathcal{O} \to 0
  \qquad
  \left\{\begin{array}{l} \psi(\lambda, \zeta) = \lambda \\ i(\xi) = (0, 2\xi) \end{array}\right.
\]
\note{à droite, un N.B. barré en diagonale : « \struck{N.B.\ On a $\lambda = \psi(x)$, $\zeta = 2x - \psi(x)T$ ; donc $x = (\lambda, \zeta)$} ».}
$T$ correspond à la section $(2, \text{\struck{\ill{}}}\ 0)$ de $\Gamma P_{2} \simeq \Gamma(2\tau + 2L)$.

\page{50}
\[
  N = (T^{2} - \delta)/4,
\]
c'est une section de $\Gamma \mathrm{Sym}^{2}(E)$, grâce à la condition $T_{0}^{2} \equiv \delta$ (4). Vérifions que la \struck{fonction} pol.
\[
  Q(x) = x^{2} - Tx + N
\]
sur $P_{1}$ est à valeurs dans $L^{\otimes 2} \subset \mathrm{Sym}^{2}(E)$. Il suffit de prouver qu'il en est ainsi de
\[
  4Q(x) = Q(2x)
\]
et on peut poser
\[
  2x = T + z
\]
on trouve
\[
  Q(T + z) = T^{2} + 2zT + z^{2} - 2T(T + z) + 4N
  \qquad\text{i.e.}
\]
\note{sous $2T(T+z)$ et $4N$ : $\psi(T+z)$ et $\psi(T+z)^{2}$.}
\[
  \boxed{Q(T + z) = \text{\struck{\ill{}}}\ z^{2} - \delta}
\]
et on a bien $z^{2} - \delta \in \Gamma L^{\otimes 2}$ \quad OK.

Description de $A$. On a
\[
  A \longrightarrow (\mathcal{O} \times \check{L}) \qquad x \longmapsto (\mathrm{Tr}(x), \varphi(x))
\]
\note{la lettre lue $\varphi$ est douteuse.}
est injectif, car le noyau est formé des $x \in \Gamma(\mathcal{O})$ ($\subset A$) tels que $\mathrm{Tr}\,x = 0$ i.e.\ $2x = 0$ i.e.\ $x = 0$. Donc
\[
  A \subset \mathcal{O} \times \check{L}
\]
On trouve
\[
  \Gamma A \simeq \bigl\{ (b, \gamma) \in \Gamma \mathcal{O} \times \Gamma \check{L} \bigm| \boxed{b_{0} = \gamma_{0} T_{0}} \bigr\}
\]
\[
  b \in \Gamma \mathcal{O}_{X}, \qquad \text{relation}\quad
  \underset{\Gamma \check{L}_{0}}{\gamma_{0}} \,.\, \underset{\Gamma L_{0}}{T_{0}} \in \Gamma \mathcal{O}_{X_{0}}
\]
\[
  \text{i.e.\ } \xi \mapsto 2\xi \qquad (b, \gamma) \longmapsto \gamma
\]
\[
  0 \to \mathcal{O}_{X} \longrightarrow A \longrightarrow \check{L} \to 0 \quad (\text{surj})
\]
Ainsi l'unité $\underline{1}$ de $A$ se décrit par $(2, 0)$.

\page{51}
\[
  \left\{\begin{array}{l}
    \text{\struck{N.B.}}\quad T(b, \gamma) = b\\
    \delta(b, \gamma) = \text{\struck{\ill{}}}\ \gamma^{2}.\delta\\[2pt]
    \text{donc}\quad N(b, \gamma) = (T^{2} - \delta)(b, \gamma)/4 = (b^{2} - \gamma^{2}\delta)/4
  \end{array}\right.
\]
C'est une section \struck{de} de $\mathcal{O}_{X}$, car
\[
  b \equiv \gamma T \ (2) \qquad\text{donc}\qquad b^{2} \equiv \gamma^{2}T^{2} \ (4)
\]
donc
\[
  b^{2} - \gamma^{2}\delta \equiv \gamma^{2}(\underbrace{T^{2} - \delta}_{\equiv 0\ (4)}) \equiv 0 \quad (4)
\]
\marginal{N.B.\ on \uncertain{en déduit} $\varphi_{N}((b,\gamma),(b',\gamma')) = \dfrac{bb' - \gamma\gamma'\delta}{2}$}

Soit $x = (b, \gamma) \in \Gamma A$, on a
\[
  x^{2} - bx + \frac{b^{2} - \gamma^{2}\delta}{4} = 0 \qquad\text{i.e.}
\]
\[
  x^{2} = bx - \frac{b^{2} - \gamma^{2}\delta}{4}
  = \Bigl( b^{2} - \frac{b^{2} - \gamma^{2}\delta}{2},\ b\gamma \Bigr)
  = \Bigl( \frac{b^{2} + \gamma^{2}\delta}{2},\ b\gamma \Bigr)
\]
Cela permet de calculer, si
\[
  x = (b, \gamma), \qquad x' = (b', \gamma')
\]
le produit $xx'$ comme
\[
  xx' = \tfrac{1}{2}\bigl( (x + x')^{2} - x^{2} - x'^{2} \bigr)
\]
i.e.
\begin{align*}
  2xx' &= \Bigl( \frac{(b + b')^{2} + (\gamma + \gamma')^{2}\delta}{2}
          - \frac{b^{2} + \gamma^{2}\delta}{2} - \frac{b'^{2} + \gamma'^{2}\delta}{2},\\
       &\qquad (b + b')(\gamma + \gamma') - b\gamma - b'\gamma' \Bigr)\\
       &= (bb' + \gamma\gamma'\delta,\ b\gamma' + b'\gamma)
\end{align*}
\marginal{mod 2 c'est bien nul car $b_{0} = \gamma_{0}T_{0}$, $b'_{0} = \gamma'_{0}T_{0}$, $T_{0}^{2} = \delta_{0}$ : $bb' + \gamma\gamma'\delta \mapsto \gamma_{0}\gamma'_{0}\delta_{0} + \gamma_{0}\gamma'_{0}\delta_{0} = 2\gamma_{0}\gamma'_{0}\delta_{0} = 0$, et $\gamma_{0}\gamma'_{0}T_{0} + \gamma_{0}\gamma'_{0}T_{0} = 2\gamma_{0}\gamma'_{0}T_{0} = 0$}
donc
\[
  \boxed{xx' = (b, \gamma)(b', \gamma') = \bigl( \tfrac{1}{2}(bb' + \gamma\gamma'\delta),\ \tfrac{1}{2}(b\gamma' + \gamma b') \bigr)}
\]

\page{52}
L'accouplement
\[
  A \times E \longrightarrow \mathcal{O} \qquad (x, f) \longmapsto \langle x, f\rangle = f(x)
\]
est donné par
\[
  \boxed{\langle (b, \gamma), (\lambda, \zeta) \rangle = \tfrac{1}{2}(\lambda b + \gamma.\zeta)}
\]
C'est OK car
\[
  \lambda_{0}\underset{\gamma_{0}T_{0}}{b_{0}} + \gamma_{0}.\underset{\lambda_{0}T_{0}}{\zeta_{0}}
  = \underbrace{\lambda_{0}\gamma_{0}T_{0}}_{} + \underbrace{\gamma_{0}\lambda_{0}T_{0}}_{}
  = 2\gamma_{0}\lambda_{0}T_{0} = 0
\]
On a
\[
  \varphi_{N}\bigl((\underset{\Gamma\mathcal{O}}{b}, \underset{\Gamma\check{L}}{\gamma}), (\underset{\Gamma\mathcal{O}}{b'}, \underset{\Gamma\check{L}}{\gamma'})\bigr)
  = \tfrac{1}{2}(bb' - \gamma\gamma'\delta)
  = \bigl\langle (b, \gamma), (b', \underset{\in \Gamma L}{-\gamma'\delta}) \bigr\rangle,
  \qquad\text{donc}
\]
$\varphi_{N} : A \to E$ est donné par
\[
  \boxed{\varphi_{N}(\underset{\Gamma\mathcal{O}}{b}, \underset{\Gamma\check{L}}{\gamma}) = (\underset{\Gamma\mathcal{O}}{b}, \underset{\Gamma L}{-\gamma\delta})}
\]
\struck{Vérifions} Utilisons la formule $Q(\varphi_{N}(x)) = +N(x)\delta$ par la formule précédente, pour $\varphi_{N}(x)$, $x = (b, \gamma)$, \struck{\ill{}}.
\[
  Q(\underset{\lambda}{b}, \underset{\zeta}{-\gamma\delta}) = - (b^{2} - \gamma^{2}\delta)/4\ \delta
  = - (\lambda^{2} + \gamma\zeta)/4\ \delta
\]
\[
  \text{\struck{\ill{}}}\ \text{donc} \qquad \gamma = -\zeta\delta^{-1}
\]
d'où
\[
  Q(\lambda, \zeta) = - (\lambda^{2} - \zeta^{2}\delta^{-1})/4\ \delta
  = \text{\struck{\ill{}}}\ (\zeta^{2} - \lambda^{2}\delta)/4
\]
i.e.
\[
  \boxed{Q(\lambda, \zeta) = (\zeta^{2} - \lambda^{2}\delta)/4}
\]
\note{les signes des deux premières lignes du calcul sont recouverts ; on les donne d'après la formule encadrée qu'il obtient.}
Il faut vérifier que
\[
  \zeta^{2} - \lambda^{2}\delta \equiv 0 \quad (4)
\]
or $\zeta \equiv \lambda T$ (2) d'où $\zeta^{2} \equiv \lambda^{2}T^{2}$ (4), donc
\[
  \zeta^{2} - \lambda^{2}\delta \equiv \lambda^{2}(\underbrace{T^{2} - \delta}_{4N}) = 4\lambda^{2}N \equiv 0 \quad (4) \qquad \text{OK}
\]
\struck{or on a $\zeta_{0} = \lambda_{0}T_{0}$,}

\struck{$\zeta_{0}^{2} - \lambda_{0}^{2}\delta_{0}$ ; $\zeta_{0}^{2} = \lambda_{0}^{2}T_{0}^{2} = \lambda_{0}^{2}\delta_{0}$ car $T_{0}^{2} = \delta_{0}$}

\page{53}
Si $y = (\lambda, \zeta) \in \Gamma P_{1}$ \add{i.e.\ $\lambda = 1$}, l'équation
\[
  Q(y) = 0 \quad\text{équivaut bien à}\quad (\zeta^{2} - \delta)/4 = 0 \quad\text{i.e.}\quad \zeta^{2} = \delta
\]
(en plus on a $\zeta_{0} = T_{0}$, sous-entendu), cela prouve que
\[
  Q'(\lambda, \zeta) = (\zeta^{2} - \lambda^{2}\delta)/4
\]
est prop.\ à $Q$, et pour voir que le facteur de prop.\ est 1, il suffit de faire \struck{$y = (\lambda, \zeta) = T$}
\[
  y = T, \qquad\text{i.e.}\qquad y = (2, 0)
\]
on a alors
\[
  \text{\struck{$Q(T) = T^{2} - 2$}}
\]
\[
  Q(y) = y^{2} - Ty\,\psi(y) + N\psi(y)^{2} = T^{2} - 2T^{2} + 4N = 4N - T^{2} = -\delta
\]
et
\[
  Q'(y) = (0 - 4\delta)/4 = -\delta \qquad \text{OK.}
\]

\subsection{Comparaison des torseurs sous $\mu_{2}$ et sous $(\mathbb{Z}/2\mathbb{Z})_{X}$}
\note{titre souligné par un trait, sur la page.}

L'hom.\ can.
\[
  (\mathbb{Z}/2\mathbb{Z})_{X} \longrightarrow \mu_{2X}
\]
définit, par torseurs associés,

\begin{tikzcd}[column sep=small, row sep=normal, nodes={font=\scriptsize}]
\mathrm{Tors}(\mathbb{Z}/2\mathbb{Z}) \arrow[r, "\mathrm{Ass}"] \arrow[d, "\wr"'] & \mathrm{Tors}(\mu_{2}) \arrow[r, "\mathrm{Oub}"] \arrow[d, "\wr"] & \mathrm{Rev\,quad}(X) \arrow[d] \\
\mathrm{Triples}(L, \delta, T_{0}) \arrow[r, "{\mathrm{Ass}'\,?}"] & \mathrm{Couples}(L, \theta) \arrow[r, "{\mathrm{Oub}'\,?}"] & \mathrm{Triples}(L, \delta, T_{0})
\end{tikzcd}

\note{sur la page, la dernière colonne est entre crochets, avec « $\otimes$-foncteur » écrit au-dessus de la flèche « Oub » ; sous la ligne du bas : à gauche « ($\delta$ inv.) », avec $\delta \in \Gamma L^{\otimes 2}$, $T_{0} \in \Gamma L_{0}$ ; au milieu « ($\theta$ inv.) », $\theta \in \Gamma L^{\otimes 2}$ ; à droite « ($\delta$ pas néc.\ inv.) », $\delta \in \Gamma L^{\otimes 2}$, $T_{0} \in \Gamma L_{0}$. Un premier signe vertical à gauche est biffé.}

Les flèches verticales sont des équiv.\ si 2 rég.\ sur $X$ (et celles du \ill{} \uncertain{aussi}, \add{\uncertain{aux conditions} \ill{} sur $X$}). Il faut déterminer les deux foncteurs horizontaux « ? ». Le premier foncteur est évidemment
\[
  \mathrm{Ass}'(L, \delta, T_{0}) = (L, \delta)
\]

\page{54}
Mais le deuxième serait plutôt
\[
  \mathrm{Oub}'(L, \theta) = (L, 4\theta, 0)
\]
Ce foncteur ne commute \add{\uncertain{mis}} pas (aux \struck{\ill{}} $\otimes$ \uncertain{à iso près}) \add{(\uncertain{si} 2 inv.\ sur $X$)}, et on en déduit que le foncteur
\[
  \mathrm{Tors}(\mu_{2X}) \longrightarrow \mathrm{Rev\,quad}(X)
\]
ne commute pas aux compositions, contrairement à :

\begin{tikzcd}
\mathrm{Tors}((\mathbb{Z}/2\mathbb{Z})_{X}) \arrow[r] \arrow[d] & \mathrm{Tors}(\mu_{2X}) \\
\mathrm{Rev\,quad}(X) &
\end{tikzcd}

\marginal{Ce n'est pas \ill{} : $(L \otimes L', 16\theta\theta')$ \ill{} $(L \otimes L', 4\theta\theta')$ \ill{} par ex.\ si $L = L' = \mathcal{O}_{X}$, $\theta = \theta' = 1$ …}
\note{note écrite verticalement dans la marge gauche, à demi lisible.}

On a en effet
\[
  (L, \delta, T_{0}) \wedge (L', \delta', T'_{0}) = (L \otimes L', \delta\delta', T_{0}T'_{0})
\]
tant dans $\underline{\mathrm{Triples\ \acute{e}tales}}\,(L, \delta, T_{0})$ que dans $\underline{\mathrm{Triples}}$ (quelconques) $(L, \delta, T_{0})$, et
\[
  (L, \theta) \wedge (L', \theta') = (L \otimes L', \theta\theta')
\]

$\underline{\mathrm{Rev\,quad}}(X)$ est une catégorie monoïdale (comm.) par un produit contracté qui prolonge celui sur les $(\mathbb{Z}/2\mathbb{Z})_{X}$-torseurs.

$\underline{\mathrm{Rev\,quad}}^{!}(X)$, sous-catégorie pleine formée des rev.\ quadratiques $X'/X$ tels que $T_{0} = 0$ i.e.\ la trace sur $X'_{0}/X_{0}$ ($X_{0} = V(2.1_{X})$) est id.\ nulle, i.e.\ $N_{X'_{0}/X_{0}}$ est additive, i.e.\ $\varphi_{N_{X'_{0}/X_{0}}}$ (forme bil.\ associée) est id.\ nulle.

\marginal{\ill{} \uncertain{qui peut être} \ill{} $\underline{\mathrm{Rev\,quad}}(X)$ \ill{} $\underline{\mathrm{Rev\,quad}}^{!}(X)$ \ill{}}
\note{note écrite en diagonale dans la marge gauche, à hauteur des deux paragraphes précédents, presque entièrement illisible.}

On définit une catégorie monoïdale mixte

\page{55}
des couples
\[
  (n, \mathfrak{X}) \qquad\text{où}\qquad
  \begin{array}{l}
    n \in \mathbb{Z}/2\mathbb{Z} \text{ est une « parité »}\\
    \mathfrak{X}/X \text{ rev.\ quadratique}\\
    \text{avec } \mathfrak{X} \in \underline{\mathrm{Rev\,quad}}^{!} \text{ si } \underline{n = 1}
  \end{array}
\]
On pose
\[
  (n, \mathfrak{X}) \wedge (n', \mathfrak{X}') =
  \left\{\begin{array}{ll}
    (n + n', \mathfrak{X} \wedge \mathfrak{X}') & \text{si $n$ ou $n'$ pair}\\
    (n + n', \mathfrak{X} \overset{!}{\wedge} \mathfrak{X}') & \text{si $n$ \emph{et} $n'$ impairs}
  \end{array}\right.
\]
où $\mathfrak{X} \overset{!}{\wedge} \mathfrak{X}'$ est un produit contracté modifié, défini si $\mathfrak{X}, \mathfrak{X}' \in \underline{\mathrm{Rev\,quad}}^{!}(X)$, et défini (si 2 rég.\ sur $X$) par la condition
\[
  (L, \delta, 0) \overset{!}{\wedge} (L', \delta', 0) = (L \otimes L', \tfrac{1}{4}\delta\delta', 0)
\]
où $\frac{1}{4}\delta\delta'$ a un sens intrinsèque, car localement on a
\[
  \begin{array}{ll}
    \delta = 4\Delta & \Delta \text{ défini mod } h \text{ avec } 4h = 0\\
    \delta' = 4\Delta' & \Delta' \text{ — } h' \text{ — } 4h' = 0
  \end{array}
\]
\[
  \delta\delta' = 16\Delta\Delta' \qquad \tfrac{1}{4}\delta\delta' \overset{\mathrm{def}}{=} 4\Delta\Delta'
\]
qui ne change pas par $\Delta \mapsto \Delta + h$ ($4h = 0$), $\Delta' \mapsto \Delta' + h'$ ($4h' = 0$).

\note{le paragraphe suivant est marqué d'un double trait vertical dans la marge gauche.}
Il faut lui donner un sens indépendamment de la condition 2 régulier sur $X$, exercice catégorique universel …

Cela \uncertain{justifiant} \emph{commutativité} \add{(à isom.\ \uncertain{can.}\ près)} et de plus, les isom.\ de comm.\ et d'\uncertain{associativité}, \ill{} les \uncertain{bilinéarités} sur $L \otimes L' \simeq L' \otimes L$, et $(L \otimes L') \otimes L'' \simeq L \otimes (L' \otimes L'')$, et \uncertain{compatibles} automatiquement avec les produits, $\delta$, $T_{0}$ etc.

\page{56}
\subsection{Rev.\ quadratiques en car.\ 2}
\note{titre souligné sur la page.}

On a :
\[
  L \rightrightarrows L^{\otimes 2} \qquad
  \left\{\begin{array}{l} u \longmapsto u^{2} \\ u \longmapsto Tu \end{array}\right.
  \qquad T \in \Gamma L \ \text{la section trace}
\]
C.N.B.\ ce sont des homs additifs, mais $u \mapsto u^{2}$ \add{\uncertain{semi-linéaire}} d'où
\[
  L \xrightarrow{\ f\ } L^{\otimes 2} : \qquad f(u) = u^{2} + Tu
\]
\[
  \check{V}(L) \longrightarrow \check{V}(L^{\otimes 2})
\]
\struck{On peut considérer $f$}

morphisme qui est étale ssi $T$ est inv., et alors $T : \mathcal{O}_{X} \xrightarrow{\ \sim\ } L$, et $f$ peut s'interpréter comme
\[
  \mathcal{O}_{X} \to \mathcal{O}_{X} : \quad u \mapsto u^{2} + u
\]

Considérons \struck{$x'$, $x''$} les deux sections « universelles » $x'$, $x''$ de $X' \subset P_{1} \subset E$, on a
\[
  x'' - x' = x'' + x' = T \qquad\text{i.e.}\qquad x'' = x' + T
\]
d'où
\[
  x''^{(2)} = x'^{(2)} + T^{2} \qquad\text{sections de}\quad P_{1} \wedge_{L} L^{2} \simeq P_{1}^{(2)}
\]

\begin{tikzcd}[column sep=normal, row sep=normal]
0 \arrow[r] & L \arrow[r] \arrow[d, "u \mapsto u^{2}"'] & E \arrow[r, "\psi"] \arrow[d, "x \mapsto x^{(2)}"] & \mathcal{O}_{X} \arrow[r] \arrow[d, "\lambda \mapsto \lambda^{2}"] & 0 \\
0 \arrow[r] & L^{\otimes 2} \arrow[r] & E^{(2)} \arrow[r, "\psi^{(2)}"] & \mathcal{O}_{X} \arrow[r] & 0
\end{tikzcd}

\[
  P_{1}^{(2)} = \psi^{(2)-1}(\{1\})
\]
\note{sous le diagramme, un bloc biffé : « \struck{$\mathrm{Sym}^{2}(E) / \{(x+y)^{(2)} - x^{(2)} - y^{(2)}\}$ \ill{}} ».}
\[
  u \in \Gamma(P_{1}^{(2)} \wedge T.P_{1}) \iff u : P_{1}^{(2)} \simeq T P_{1} \quad \text{iso de $L^{\otimes 2}$-torseurs}
\]
\[
  \iff P_{1} \xrightarrow{\ u''\ } P_{1}^{(2)} \quad \text{compatible avec}\quad L \to L^{\otimes 2},\ u \mapsto T.u
\]

\emph{N.B.} $f : \check{V}(L) \to \check{V}(L^{\otimes 2})$ est fid.\ plat (car \ill{}) et son noyau est un sous-groupe \add{\uncertain{fini}} quadratique de $\check{V}(L)$, noté $G_{T}$, étale ssi $T$ est inv.\ (\ill{} $V(T)$, schéma des zéros de $T$).

Soit $x \in \Gamma(P_{1})$, i.e.\ $x \in \Gamma(P_{1})$, \uncertain{tel que}
\[
  Q(x) = 0 \qquad x^{2} + Tx + N = 0
\]
Écrivons
\[
  Q(x + u) = 0 \quad (u \in \Gamma L) \qquad x^{2} + u^{2} + Tx + Tu + N = 0, \quad\text{i.e.}\quad u^{2} + Tu = 0
\]

\page{57}
Donc $X'/X$ est un \emph{torseur} sous $G_{T}$ ! \ill{}

\emph{Th.} On a une équivalence de catégories \add{(groupoïdes)} fibrées sur $\mathrm{Sch}/\mathbb{F}_{2}$ (plus gén.\ sur les topos loc.\ annelés en car.\ 2), entre la catégorie des rev.\ quadratiques relatifs (sur $X \in \mathrm{Ob}\,\mathrm{Sch}/\mathbb{F}_{2}$), et les données (sur un tel $X$)
\[
  (L, \overset{P_{1}}{\text{\struck{$\ast$}}}, T, \alpha), \quad\text{où}\quad
  \left\{\begin{array}{l}
    L \text{ module inv.\ sur } X, \ P_{1} \ L\text{-tors.}\\
    \quad \text{(i.e.\ $P_{1}$ un fibré en droites affine sur $X$)}\\
    T \text{ section de } L \text{ (module des translations de $P_{1}$)}\\
    \alpha \text{ section de } P_{1} \wedge_{L} (L^{\otimes 2}, f_{T} : L \to L^{\otimes 2})
  \end{array}\right.
\]
où
\[
  f_{T} : L \to L^{\otimes 2} \quad\text{défini par}\quad f_{T}(u) = u^{2} + T.u
\]
(c'est un \struck{morphisme} hom.\ additif, \uncertain{semi-linéaire}, $W(L) \to W(L^{\otimes 2})$, dont le noyau est un groupe \add{\uncertain{fini}} quadratique $G_{T}$ sur $X$).

Au système $\Sigma = (L, P_{1}, T, \alpha)$ est associé le $X$-schéma
\[
  X'_{\Sigma} = f_{T, P_{1}}^{-1}(V(\alpha))
\]
(torseur sous $G_{T}$), qui est \uncertain{quadratique} sur $X$. Plus de doute, c'est lui !

\marginal{N.B.\ $u : P_{1} \to P_{1}^{(2)}$ compatible avec $u \mapsto Tu$ est donné par $u : x \mapsto Tx + N$, on a alors ($Q(x) = x^{2} + Tx + N$), $u(x) = x^{2} + Q(x)$ dans \ill{}, \struck{\ill{}} l'équation $Q(x) = 0$ qui \uncertain{définit} $X'/X$ équivaut à : $\boxed{u(x) = x^{2}}$, $P_{1} \rightrightarrows P_{1}^{(2)}$ ($u$ et $x \mapsto x^{2}$).}
\note{note encadrée d'un trait, à droite du texte.}

Comment s'explicite le produit contracté ?
\[
  (L, T, P_{1}, \alpha) \wedge (L', T', P'_{1}, \alpha') \overset{?}{=} (L \otimes L', TT', ?, ?)
\]
\struck{Ainsi} Ainsi :
\[
  \begin{array}{lll}
    G_{T} \xrightarrow{\ d_{T'}\ } G_{TT'} & \text{induit par} & d_{T'} : L \to L \otimes L', \ u \mapsto uT'\\
    G_{T'} \longrightarrow G_{TT'} & \text{—} & d_{T} : L' \to L \otimes L', \ u' \mapsto Tu'
  \end{array}
\]
d'où
\[
  G_{T} \times G_{T'} \longrightarrow G_{TT'} \quad\text{hom.\ de groupes}
\]
et si \struck{$\mathfrak{X}$} $\mathfrak{X}$ torseur sous $G_{T}$, $\mathfrak{X}'$ sous $G_{T'}$, on en déduit par \uncertain{extension} du groupe un torseur $\mathfrak{X} \wedge \mathfrak{X}'$ sous $G_{TT'}$ !

\page{58}
\struck{Trois} \add{Quatre} catégories fibrées sur $(\mathrm{Sch})$ (plus gén.\ sur la catégorie des topos loc.\ annelés …), canoniquement équivalentes
\begin{enumerate}
  \item[①] Catégorie des rev.\ quadratiques $X'/X$ de $X \in \mathrm{Ob}\,\mathrm{Sch}$,
  \item[①'] Catégorie des alg.\ \add{$V(A)$} commutatives loc.\ libres de rang 2 sur $\mathcal{O}_{X}$
  \item[②] Catégorie des triples $(A, N, e)$, où $(A, N)$ est un module quadratique loc.\ libre de rang 2 sur $\mathcal{O}_{X}$, et où $e$ est une section de $A$ telle que
  \[
    N(e) = 1
  \]
  \item[③] Catégorie des triples $(L, P, Q)$, où $L$ est un $\mathcal{O}_{X}$-Module inv.\ \add{(identifié au schéma en droites homogène $W(L)$)}, $P$ un $L$-torseur, et $Q$ une application polynomiale
  \[
    P \longrightarrow L^{\otimes 2}
  \]
  dont la partie homogène de degré 2 soit $u \mapsto u^{2}$, i.e.\ on suppose que
  \[
    Q(x + u) = Q(x) + u^{2} + \mathrm{Lin}_{u}(x)
  \]
  où
  \[
    \mathrm{Lin}_{u} : P \to L_{S}^{\otimes 2}
  \]
  est une application \emph{affine} \add{(\uncertain{applis.\ dessus})} (quand $u$ est une section loc.\ de $L$ sur $S/X$)
  \item[③'] Catégorie des triples ( \struck{$0 \to L$} $0 \to L \xrightarrow{\ i\ } E \xrightarrow{\ \psi\ } \mathcal{O}_{X} \to 0, Q$) d'une suite exacte $0 \to L \to E \to \mathcal{O}_{X} \to 0$, avec $L$ inv., et d'une forme quadratique $Q$ sur $E$ à val.\ dans $L^{\otimes 2}$, dont la restriction à $L$ soit $u \mapsto u^{\otimes 2}$.
\end{enumerate}
\note{dans la marge gauche, des flèches relient ① et ①' (double flèche), ①' et ② (flèche vers le bas), ③ et ① (grande flèche courbe), ③ et ③' (double flèche).}

On va transformer ③', en \struck{\ill{}} considérant $L^{\otimes 2}$ comme plongé dans $\mathrm{Sym}^{2}(E)$ ($\simeq \mathrm{Quad}(A)$, où $A \overset{\mathrm{df}}{=} \check{E}$), laquelle a une suite de composition

\page{59}
\[
  0 \subset \underbrace{L^{\otimes 2}}_{L^{\otimes 2}} \subset \underbrace{L \otimes E}_{L} \subset \underbrace{\mathrm{Sym}^{2}(E)}_{\mathcal{O}_{X}}
\]
\note{sous chaque inclusion, l'accolade donne le quotient successif.}
Considérons \struck{écrivons}
\[
  Q(x) - x^{2} \ : \ P \to \mathrm{Sym}^{2}(E)
\]
fonction polynomiale de degré \struck{\ill{}} $\leq 2$, \struck{\ill{}} dont la partie \struck{\ill{} homogène} \add{quadratique} \struck{$Q(x + u) -$}
\[
  L \to L^{\otimes 2}
\]
est nulle. Donc c'est une application \emph{affine}.

\struck{i.e.\ $Q(x) = x^{2} + \mathrm{Af}(x)$, $\mathrm{Af} : P \to \mathrm{Sym}^{2}E$ application affine,}

\struck{dont la partie linéaire est un hom.\ $L \to \mathrm{Sym}^{2}(E)$ défini par un élément $T$ de $\check{L} \otimes \mathrm{Sym}^{2}E \supset E$}
\note{ces deux lignes sont barrées de longs traits obliques.}

Utilisant une description de $N$ et $E$ par une section $e$ de $P_{1}$, i.e.\ $e \in \Gamma E$, $\psi(e) = 1$, on aura, posant
\[
  \left\{\begin{array}{l} x = e + u, \\ u = x - e \end{array}\right.
  \qquad u \in \Gamma L,
\]
\add{et écrivons \ill{}}
\[
  Q(e + u) = u^{2} + bu + c
  \qquad
  \begin{array}{l} u \in \Gamma L\\ b \in \Gamma L\\ c \in \Gamma L^{\otimes 2} \end{array}
\]
\begin{align*}
  Q(\underbrace{e + u}_{x}) - (\underbrace{e + u}_{x})^{2}
  &= bu + c - 2eu - e^{2} = \underset{x - e}{u}(b - 2e) + c - e^{2}\\
  &= x(b - 2e) \ \text{\struck{\ill{}}} + (e^{2} + c - be)
\end{align*}
i.e.
\[
  \left\{\begin{array}{l}
    Q(x) = x^{2} - xT + N\\[2pt]
    \text{où}\quad \boxed{T = 2e - b \in E, \quad N = e^{2} + c - be \in \mathrm{Sym}^{2}(E)}\\[2pt]
    \text{avec}\quad \boxed{\psi(T) = 2, \ \varepsilon(N) = 1}
  \end{array}\right.
\]
\note{sous $N$ : $e^{2} \in \Gamma L^{2}$ \uncertain{(sic)}, $c \in \Gamma L^{\otimes 2}$, $be \in L \otimes E \subset \mathrm{Sym}^{2}E$ ; à droite, $\varepsilon : \mathrm{Sym}^{2}(E) \to \mathcal{O}_{X} \to 0$, avec $0 \to L \otimes E \to$.}

\struck{\ill{}} On vérifie que $T$, donc aussi $N$, ne dépend pas du scindage choisi. [$T$ est élément d'un torseur \add{$P_{2} = \psi^{-1}(2)$} sous $L$, \struck{\ill{}} et $N$ (pour $T$ choisi), avec la condition que $N$ soit à valeurs dans $L^{\otimes 2}$ \add{(\uncertain{i.e.\ le choix de} $L^{\otimes 2}$)} est élément d'un torseur sous $L^{\otimes 2} \subset \mathrm{Sym}^{2}(E)$, donc \struck{$Q$} \add{\ill{}} $Q \leftrightarrow (T, N)$ est \uncertain{celui}

\page{60}
d'un élément d'un torseur sous $L^{\otimes 2} \otimes \check{E}$ \add{$= L^{\otimes 2} \otimes A$}, \struck{\ill{}} qui peut être considéré comme l'espace des applications affines de $P_{1}$ dans $L^{\otimes 2}$, \add{\uncertain{sur} $\check{E} \otimes L^{\otimes 2}$ des} des applications linéaires de $E$ dans $L^{\otimes 2}$, \struck{\ill{}} c'est-à-dire
\[
  0 \to L^{\otimes 2} \to L^{\otimes 2} \otimes \check{E} \to L \to 0
\]
(car $0 \to \mathcal{O}_{X} \to \check{E} \to \check{L} \to 0$ et on multiplie par $L^{\otimes 2}$)
\[
  L^{\otimes 2} \otimes \check{E} \simeq L \otimes E
\]
\note{l'isomorphisme est écrit verticalement, marqué « can » : « (pour tout module loc.\ libre de rg 2, $E$, on a $L \simeq \det E$) ».}
\[
  \text{i.e.}\qquad E \simeq \check{E} \otimes L \ (\simeq \mathrm{Hom}(E, L)), \qquad L = \det E
\]

L'application affine
\[
  P_{1} \longrightarrow L^{\otimes 2}
\]
définie par un
\[
  \xi \in L.E \simeq L \otimes E \subset \mathrm{Sym}^{2}E
\]
est définie par
\[
  x \longmapsto x \wedge \xi
  \qquad\qquad
  \begin{array}{l}
    E \otimes E \to L, \quad (x \otimes y) \mapsto x \wedge y\\
    L \otimes E \otimes E \to L \otimes L = L^{\otimes 2}\\
    (\xi, x) \mapsto \xi \wedge x = -x \wedge \xi
  \end{array}
\]
\note{le second membre $x \wedge \xi$ est surchargé et en partie biffé.}

… ]

Inversement, si la \uncertain{suite exacte} $E$ de $\mathcal{O}_{X}$ par $L$ étant donnée, (sans $Q$ encore) si on se donne
\[
  T \in E, \qquad N \in \mathrm{Sym}^{2}(E),
\]
quand correspondent-ils à une $Q$ comme dans 3') ? Il faut, \uncertain{supposons-les},
\[
  \varphi_{1}(T) = 2, \qquad \varphi_{2}(N) = 1
\]
(i.e.\ en termes de la « section unité » $e$ de $A = \check{E}$ : $T(e) = 2$, $N(e) = 1$) plus une condition qui équivaudrait à
\[
  N(x + \lambda e) = N(x) + \lambda T(x) + \lambda^{2}
  \qquad\text{(on aura $N(x + \lambda e) = N(x) + \lambda T(x) + \lambda^{2}$)}
\]
i.e.
\[
  T(x) = \varphi_{N}(x, e)
\]
En tous cas on doit avoir : \emph{$N$ détermine $T$}
\note{l'argument se poursuit au-delà de ce lot.}

\end{document}
